\IfFileExists{stacks-project.cls}{%
\documentclass{stacks-project}
}{%
\documentclass{amsart}
}

% For dealing with references we use the comment environment
\usepackage{verbatim}
\newenvironment{reference}{\comment}{\endcomment}
%\newenvironment{reference}{}{}
\newenvironment{slogan}{\comment}{\endcomment}
\newenvironment{history}{\comment}{\endcomment}

% For commutative diagrams we use Xy-pic
\usepackage[all]{xy}

% We use 2cell for 2-commutative diagrams.
\xyoption{2cell}
\UseAllTwocells

% We use multicol for the list of chapters between chapters
\usepackage{multicol}

% This is generally recommended for better output
\usepackage{lmodern}
\usepackage[T1]{fontenc}

% For cross-file-references
\usepackage{xr-hyper}

% Package for hypertext links:
\usepackage{hyperref}

% For any local file, say "hello.tex" you want to link to please
% use \externaldocument[hello-]{hello}
\externaldocument[introduction-]{introduction}
\externaldocument[conventions-]{conventions}
\externaldocument[sets-]{sets}
\externaldocument[categories-]{categories}
\externaldocument[topology-]{topology}
\externaldocument[sheaves-]{sheaves}
\externaldocument[sites-]{sites}
\externaldocument[stacks-]{stacks}
\externaldocument[fields-]{fields}
\externaldocument[algebra-]{algebra}
\externaldocument[brauer-]{brauer}
\externaldocument[homology-]{homology}
\externaldocument[derived-]{derived}
\externaldocument[simplicial-]{simplicial}
\externaldocument[more-algebra-]{more-algebra}
\externaldocument[smoothing-]{smoothing}
\externaldocument[modules-]{modules}
\externaldocument[sites-modules-]{sites-modules}
\externaldocument[injectives-]{injectives}
\externaldocument[cohomology-]{cohomology}
\externaldocument[sites-cohomology-]{sites-cohomology}
\externaldocument[dga-]{dga}
\externaldocument[dpa-]{dpa}
\externaldocument[sdga-]{sdga}
\externaldocument[hypercovering-]{hypercovering}
\externaldocument[schemes-]{schemes}
\externaldocument[constructions-]{constructions}
\externaldocument[properties-]{properties}
\externaldocument[morphisms-]{morphisms}
\externaldocument[coherent-]{coherent}
\externaldocument[divisors-]{divisors}
\externaldocument[limits-]{limits}
\externaldocument[varieties-]{varieties}
\externaldocument[topologies-]{topologies}
\externaldocument[descent-]{descent}
\externaldocument[perfect-]{perfect}
\externaldocument[more-morphisms-]{more-morphisms}
\externaldocument[flat-]{flat}
\externaldocument[groupoids-]{groupoids}
\externaldocument[more-groupoids-]{more-groupoids}
\externaldocument[etale-]{etale}
\externaldocument[chow-]{chow}
\externaldocument[intersection-]{intersection}
\externaldocument[pic-]{pic}
\externaldocument[weil-]{weil}
\externaldocument[adequate-]{adequate}
\externaldocument[dualizing-]{dualizing}
\externaldocument[duality-]{duality}
\externaldocument[discriminant-]{discriminant}
\externaldocument[derham-]{derham}
\externaldocument[local-cohomology-]{local-cohomology}
\externaldocument[algebraization-]{algebraization}
\externaldocument[curves-]{curves}
\externaldocument[resolve-]{resolve}
\externaldocument[models-]{models}
\externaldocument[functors-]{functors}
\externaldocument[equiv-]{equiv}
\externaldocument[pione-]{pione}
\externaldocument[etale-cohomology-]{etale-cohomology}
\externaldocument[proetale-]{proetale}
\externaldocument[relative-cycles-]{relative-cycles}
\externaldocument[more-etale-]{more-etale}
\externaldocument[trace-]{trace}
\externaldocument[crystalline-]{crystalline}
\externaldocument[spaces-]{spaces}
\externaldocument[spaces-properties-]{spaces-properties}
\externaldocument[spaces-morphisms-]{spaces-morphisms}
\externaldocument[decent-spaces-]{decent-spaces}
\externaldocument[spaces-cohomology-]{spaces-cohomology}
\externaldocument[spaces-limits-]{spaces-limits}
\externaldocument[spaces-divisors-]{spaces-divisors}
\externaldocument[spaces-over-fields-]{spaces-over-fields}
\externaldocument[spaces-topologies-]{spaces-topologies}
\externaldocument[spaces-descent-]{spaces-descent}
\externaldocument[spaces-perfect-]{spaces-perfect}
\externaldocument[spaces-more-morphisms-]{spaces-more-morphisms}
\externaldocument[spaces-flat-]{spaces-flat}
\externaldocument[spaces-groupoids-]{spaces-groupoids}
\externaldocument[spaces-more-groupoids-]{spaces-more-groupoids}
\externaldocument[bootstrap-]{bootstrap}
\externaldocument[spaces-pushouts-]{spaces-pushouts}
\externaldocument[spaces-chow-]{spaces-chow}
\externaldocument[groupoids-quotients-]{groupoids-quotients}
\externaldocument[spaces-more-cohomology-]{spaces-more-cohomology}
\externaldocument[spaces-simplicial-]{spaces-simplicial}
\externaldocument[spaces-duality-]{spaces-duality}
\externaldocument[formal-spaces-]{formal-spaces}
\externaldocument[restricted-]{restricted}
\externaldocument[spaces-resolve-]{spaces-resolve}
\externaldocument[formal-defos-]{formal-defos}
\externaldocument[defos-]{defos}
\externaldocument[cotangent-]{cotangent}
\externaldocument[examples-defos-]{examples-defos}
\externaldocument[algebraic-]{algebraic}
\externaldocument[examples-stacks-]{examples-stacks}
\externaldocument[stacks-sheaves-]{stacks-sheaves}
\externaldocument[criteria-]{criteria}
\externaldocument[artin-]{artin}
\externaldocument[quot-]{quot}
\externaldocument[stacks-properties-]{stacks-properties}
\externaldocument[stacks-morphisms-]{stacks-morphisms}
\externaldocument[stacks-limits-]{stacks-limits}
\externaldocument[stacks-cohomology-]{stacks-cohomology}
\externaldocument[stacks-perfect-]{stacks-perfect}
\externaldocument[stacks-introduction-]{stacks-introduction}
\externaldocument[stacks-more-morphisms-]{stacks-more-morphisms}
\externaldocument[stacks-geometry-]{stacks-geometry}
\externaldocument[moduli-]{moduli}
\externaldocument[moduli-curves-]{moduli-curves}
\externaldocument[examples-]{examples}
\externaldocument[exercises-]{exercises}
\externaldocument[guide-]{guide}
\externaldocument[desirables-]{desirables}
\externaldocument[coding-]{coding}
\externaldocument[obsolete-]{obsolete}
\externaldocument[fdl-]{fdl}
\externaldocument[index-]{index}

% Theorem environments.
%
\theoremstyle{plain}
\newtheorem{theorem}[subsection]{Theorem}
\newtheorem{proposition}[subsection]{Proposition}
\newtheorem{lemma}[subsection]{Lemma}

\theoremstyle{definition}
\newtheorem{definition}[subsection]{Definition}
\newtheorem{example}[subsection]{Example}
\newtheorem{exercise}[subsection]{Exercise}
\newtheorem{situation}[subsection]{Situation}

\theoremstyle{remark}
\newtheorem{remark}[subsection]{Remark}
\newtheorem{remarks}[subsection]{Remarks}

\numberwithin{equation}{subsection}

% Macros
%
\def\lim{\mathop{\mathrm{lim}}\nolimits}
\def\colim{\mathop{\mathrm{colim}}\nolimits}
\def\Spec{\mathop{\mathrm{Spec}}}
\def\Hom{\mathop{\mathrm{Hom}}\nolimits}
\def\Ext{\mathop{\mathrm{Ext}}\nolimits}
\def\SheafHom{\mathop{\mathcal{H}\!\mathit{om}}\nolimits}
\def\SheafExt{\mathop{\mathcal{E}\!\mathit{xt}}\nolimits}
\def\Sch{\mathit{Sch}}
\def\Mor{\mathop{\mathrm{Mor}}\nolimits}
\def\Ob{\mathop{\mathrm{Ob}}\nolimits}
\def\Sh{\mathop{\mathit{Sh}}\nolimits}
\def\NL{\mathop{N\!L}\nolimits}
\def\CH{\mathop{\mathrm{CH}}\nolimits}
\def\proetale{{pro\text{-}\acute{e}tale}}
\def\etale{{\acute{e}tale}}
\def\QCoh{\mathit{QCoh}}
\def\Ker{\mathop{\mathrm{Ker}}}
\def\Im{\mathop{\mathrm{Im}}}
\def\Coker{\mathop{\mathrm{Coker}}}
\def\Coim{\mathop{\mathrm{Coim}}}

% Boxtimes
%
\DeclareMathSymbol{\boxtimes}{\mathbin}{AMSa}{"02}

%
% Macros for moduli stacks/spaces
%
\def\QCohstack{\mathcal{QC}\!\mathit{oh}}
\def\Cohstack{\mathcal{C}\!\mathit{oh}}
\def\Spacesstack{\mathcal{S}\!\mathit{paces}}
\def\Quotfunctor{\mathrm{Quot}}
\def\Hilbfunctor{\mathrm{Hilb}}
\def\Curvesstack{\mathcal{C}\!\mathit{urves}}
\def\Polarizedstack{\mathcal{P}\!\mathit{olarized}}
\def\Complexesstack{\mathcal{C}\!\mathit{omplexes}}
% \Pic is the operator that assigns to X its picard group, usage \Pic(X)
% \Picardstack_{X/B} denotes the Picard stack of X over B
% \Picardfunctor_{X/B} denotes the Picard functor of X over B
\def\Pic{\mathop{\mathrm{Pic}}\nolimits}
\def\Picardstack{\mathcal{P}\!\mathit{ic}}
\def\Picardfunctor{\mathrm{Pic}}
\def\Deformationcategory{\mathcal{D}\!\mathit{ef}}
\usepackage{cleveref}
\crefname{originaldefinition}{original definition}{original definitions}
\Crefname{originaldefinition}{Original definition}{Original definitions}
\makeatletter
\newlabel{situation-q}{{082M}{source}{}{sourcepointer-1}{}}
\newlabel{situation-q@cref}{{[originaldefinition][0][]082M}{[1][0][]source}{}{}{}}
\makeatother
\crefname{originallemma}{original lemma}{original lemmas}
\Crefname{originallemma}{Original lemma}{Original lemmas}
\makeatletter
\newlabel{more-morphisms-lemma-pushout-along-thickening}{{07RT}{source}{}{sourcepointer-2}{}}
\newlabel{more-morphisms-lemma-pushout-along-thickening@cref}{{[originallemma][0][]07RT}{[1][0][]source}{}{}{}}
\makeatother
\crefname{originallemma}{original lemma}{original lemmas}
\Crefname{originallemma}{Original lemma}{Original lemmas}
\makeatletter
\newlabel{spaces-pushouts-lemma-space-over-pushout-flat-modules}{{08KV}{source}{}{sourcepointer-3}{}}
\newlabel{spaces-pushouts-lemma-space-over-pushout-flat-modules@cref}{{[originallemma][0][]08KV}{[1][0][]source}{}{}{}}
\makeatother
\crefname{originallemma}{original lemma}{original lemmas}
\Crefname{originallemma}{Original lemma}{Original lemmas}
\makeatletter
\newlabel{more-algebra-lemma-module-over-fibre-product-bis}{{08IG}{source}{}{sourcepointer-4}{}}
\newlabel{more-algebra-lemma-module-over-fibre-product-bis@cref}{{[originallemma][0][]08IG}{[1][0][]source}{}{}{}}
\makeatother
\crefname{originallemma}{original lemma}{original lemmas}
\Crefname{originallemma}{Original lemma}{Original lemmas}
\makeatletter
\newlabel{more-algebra-lemma-surjection-module-over-fibre-product}{{08KJ}{source}{}{sourcepointer-5}{}}
\newlabel{more-algebra-lemma-surjection-module-over-fibre-product@cref}{{[originallemma][0][]08KJ}{[1][0][]source}{}{}{}}
\makeatother
\crefname{originalremark}{original remark}{original remarks}
\Crefname{originalremark}{Original remark}{Original remarks}
\makeatletter
\newlabel{more-algebra-remark-relative-modules-over-fibre-product}{{08KR}{source}{}{sourcepointer-6}{}}
\newlabel{more-algebra-remark-relative-modules-over-fibre-product@cref}{{[originalremark][0][]08KR}{[1][0][]source}{}{}{}}
\makeatother
\title{The quotient functor satisfies the Rim-Schlessinger condition on pushouts}
\author{The Stacks Project Authors}
\date{stacks-project:quot.tex:local-snapshot-20261012; GNU FDL 1.2 or later}
\begin{document}
\maketitle
\setcounter{subsection}{1}
\noindent\textit{Editorial note.} All delivered mathematical text is the author's own, copied from the pinned Stacks Project chapter file quot.tex; no statement or proof step was rewritten, reordered or removed. The counted span is the author environment and its own proof environment. References inside the span are handed to the assembler as navigation stubs that carry the canonical tag of the cited tag only; whether each of them is load-bearing is recorded on the stub itself.\par\medskip
\noindent\textit{Allowed standard prerequisites (editorial).}\begin{itemize}
\item Definition situation-q (canonical tag 082M). cited at statement line 2215. Standing situation.
\item Lemma more-morphisms-lemma-pushout-along-thickening (canonical tag 07RT). cited at proof line 2224. The pushout square.
\item Lemma spaces-pushouts-lemma-space-over-pushout-flat-modules (canonical tag 08KV). cited at proof line 2241. Construction of A'.
\item Lemma more-algebra-lemma-module-over-fibre-product-bis (canonical tag 08IG). cited at proof line 2262. First arrow surjective.
\item Lemma more-algebra-lemma-surjection-module-over-fibre-product (canonical tag 08KJ). cited at proof lines 2263-2265. Second arrow surjective.
\item Remark more-algebra-remark-relative-modules-over-fibre-product (canonical tag 08KR). cited at proof line 2272. Finitely presented case.
\end{itemize}\medskip
\ifdefined\hypertarget\hypertarget{sourcepointer-1}{}\fi
\noindent\textit{Source reference.} Original definition 082M: sources/stacks-chow/tags.txt, source lines 10257--10257. Cited at statement line 2215: the lemma is phrased in this situation.\par\medskip
\ifdefined\hypertarget\hypertarget{sourcepointer-2}{}\fi
\noindent\textit{Source reference.} Original lemma 07RT: sources/stacks-chow/tags.txt, source lines 9878--9878. Cited at proof line 2224: the pushout square exists.\par\medskip
\ifdefined\hypertarget\hypertarget{sourcepointer-3}{}\fi
\noindent\textit{Source reference.} Original lemma 08KV: sources/stacks-chow/tags.txt, source lines 10898--10898. Cited at proof line 2241: produces A' flat over Y'.\par\medskip
\ifdefined\hypertarget\hypertarget{sourcepointer-4}{}\fi
\noindent\textit{Source reference.} Original lemma 08IG: sources/stacks-chow/tags.txt, source lines 10814--10814. Cited at proof line 2262: first arrow surjective.\par\medskip
\ifdefined\hypertarget\hypertarget{sourcepointer-5}{}\fi
\noindent\textit{Source reference.} Original lemma 08KJ: sources/stacks-chow/tags.txt, source lines 10887--10887. Cited at proof lines 2263-2265: second arrow surjective.\par\medskip
\ifdefined\hypertarget\hypertarget{sourcepointer-6}{}\fi
\noindent\textit{Source reference.} Original remark 08KR: sources/stacks-chow/tags.txt, source lines 10894--10894. Cited at proof line 2272: the Q\textasciicircum{}fp case in the algebra setting.\par\medskip
\begin{lemma}
\label{lemma-q-RS-star}
In Situation \ref{situation-q}. Let
$$
\xymatrix{
Z \ar[r] \ar[d] & Z' \ar[d] \\
Y \ar[r] & Y'
}
$$
be a pushout in the category of schemes over $B$ where
$Z \to Z'$ is a thickening and $Z \to Y$ is affine, see
More on Morphisms, Lemma \ref{more-morphisms-lemma-pushout-along-thickening}.
Then the natural map
$$
Q_{\mathcal{F}/X/B}(Y') \longrightarrow
Q_{\mathcal{F}/X/B}(Y) \times_{Q_{\mathcal{F}/X/B}(Z)} Q_{\mathcal{F}/X/B}(Z')
$$
is bijective. If $X \to B$ is locally of finite presentation, then
the same thing is true for $Q^{fp}_{\mathcal{F}/X/B}$.
\end{lemma}

\begin{proof}
Let us construct an inverse map. Namely, suppose we have
$\mathcal{F}_Y \to \mathcal{A}$,
$\mathcal{F}_{Z'} \to \mathcal{B}'$, and an isomorphism
$\mathcal{A}|_{X_Z} \to \mathcal{B}'|_{X_Z}$
compatible with the given surjections.
Then we apply Pushouts of Spaces, Lemma
\ref{spaces-pushouts-lemma-space-over-pushout-flat-modules}
to get a quasi-coherent module $\mathcal{A}'$ on $X_{Y'}$
flat over $Y'$. Since this sheaf is constructed as a fibre product
(see proof of cited lemma) there is a canonical map
$\mathcal{F}_{Y'} \to \mathcal{A}'$. That this map is surjective
can be seen because it factors as
$$
\begin{matrix}
\mathcal{F}_{Y'} \\
\downarrow \\
(X_Y \to X_{Y'})_*\mathcal{F}_Y
\times_{(X_Z \to X_{Y'})_*\mathcal{F}_Z}
(X_{Z'} \to X_{Y'})_*\mathcal{F}_{Z'} \\
\downarrow \\
\mathcal{A}' =
(X_Y \to X_{Y'})_*\mathcal{A}
\times_{(X_Z \to X_{Y'})_*\mathcal{A}|_{X_Z}}
(X_{Z'} \to X_{Y'})_*\mathcal{B}'
\end{matrix}
$$
and the first arrow is surjective by
More on Algebra, Lemma \ref{more-algebra-lemma-module-over-fibre-product-bis}
and the second by
More on Algebra, Lemma
\ref{more-algebra-lemma-surjection-module-over-fibre-product}.

\medskip\noindent
In the case of $Q^{fp}_{\mathcal{F}/X/B}$ all we have to show is that
the construction above produces a finitely presented module.
This is explained in
More on Algebra, Remark
\ref{more-algebra-remark-relative-modules-over-fibre-product}
in the commutative algebra setting. The current case of modules
over algebraic spaces follows from this
by \'etale localization.
\end{proof}

\end{document}
